\documentclass[journal,twoside,web]{ieeecolor}

\usepackage{generic}
\usepackage{cite}
\usepackage{amsfonts, amsmath, amssymb, bbm}
\usepackage{color}
\usepackage{comment}
\usepackage{import}
\usepackage{mathtools}
\usepackage[short]{optidef}
\usepackage{graphicx}
%\usepackage{ulem}
\usepackage{hyperref}
\usepackage{algorithm}


%---------

%\documentclass[letterpaper, 10pt, conference]{ieeeconf}  % Comment this line out if you need a4paper

%\documentclass[a4paper, 10pt, conference]{ieeeconf}      % Use this line for a4 paper

%\IEEEoverridecommandlockouts                              % This command is only needed if
                                                          % you want to use the \thanks command

\newtheorem{definition}{Definition}
\newtheorem{theorem}{Theorem}
\newtheorem{corollary}{Corollary}
\newtheorem{lemma}{Lemma}
\newtheorem{assumption}{Assumption}
\newtheorem{proposition}{Proposition}
\newtheorem{example}{Example}
\newtheorem{remark}{Remark}
\newtheorem{problem}{Problem}
\newtheorem{observation}{Observation}



\newcommand{\state}{\mathbf{x}}
\newcommand{\disturbance}{\mathbf{\widehat w}}
\newcommand{\operator}{\mathcal{M}}
\newcommand{\magnitude}{\mathcal{D}}

\pagestyle{empty} % Removes all the page numbers (except for the title page)


\begin{document}

\title{\LARGE \bf Context-Enriched Performance Boosting via Operator Decomposition}


\author{Leonardo Massai, Sebastiano Messina, Nicolas Kirsch, and Giancarlo Ferrari-Trecate % <-this % stops a space
	%\thanks{This work was not supported by any organization}% <-this % stops a space
	\thanks{This research has been supported by the Swiss National Science Foundation under the NCCR Automation (grant agreement 51NF40\_180545) and the NECON project (grant number 200021-219431).}
    \thanks{The authors are with the Institute of Mechanical Engineering, Ecole Polytechnique Fédérale de Lausanne (EPFL), CH-1015 Lausanne, Switzerland. (email:}
}


\maketitle
\thispagestyle{empty}
 

%%%%%%%%%%%%%%%%%%%%%%%%%%%%%%%%%%%%%%%%%%%%%%%%%%%%%%%%%%%%%%%%%%%%%%%%%%%%%%%%
\begin{abstract}
Performance Boosting (PB) is a control framework that, for a pre-stabilized system subject to \(\mathcal L_p\) process disturbances, parametrizes all and only controllers that preserve closed-loop \(\mathcal L_p\)-stability. This parametrization is based on a causal \(\mathcal L_p\)-stable operator mapping the reconstructed disturbance signal to the corresponding corrective control action. Although PB allows optimization over the full class of stability-preserving controllers, searching this space efficiently can be difficult. We show that enriching the PB operator with contextual information can simplify learning and dramatically improve performance, particularly in multi-input settings where context helps allocate corrective actions across control channels. To this end, we propose a structured decomposition for context-enriched \(\mathcal L_p\)-stable operators combining a stable disturbance-processing dynamical module with a uniformly bounded context-dependent mixer. This yields a modular and interpretable architecture compatible with PB stability analysis. We further prove that, on a broad admissible disturbance regime characterized by a uniform envelope condition, the proposed factorization is necessary and sufficient for the corresponding class of causal extended operators, and is therefore non-conservative on that regime. This setting naturally covers many transient disturbances relevant in control and is compatible with fading-memory dynamical models such as stable state-space architectures. A numerical study illustrates the benefits of the proposed structure in a representative PB scenario.
\end{abstract}



\begin{IEEEkeywords}
Performance Boosting control, stability-preserving parametrization, contextual control, operator decomposition, $\mathcal L_p$ stability, multi-input control.
\end{IEEEkeywords}
\section{Introduction}

Modern control systems face increasing complexity in both their dynamics and objectives. In many modern applications, performance objectives depend explicitly on exogenous signals such as dynamic energy prices for battery management \cite{ali2025optimal} or real-time obstacle positions in mobile robotics \cite{10680398}.Traditional controllers, however, remain largely agnostic to these contextual inputs, limiting their effectiveness for sophisticated tasks requiring environmental awareness or anticipatory adaptation.

The rapid expansion of sensor networks and data availability further emphasizes the need for control architectures that explicitly leverage exogenous contextual information. Improving control performance through contextual awareness has therefore become an active research direction. By directly feeding exogenous signals to the controller, the system can better adapt to its operating environment and achieve more complex objectives. 

Several works have investigated this idea within Model Predictive Control (MPC). In \cite{10680398}, contextual information such as the positions and velocities of humans and obstacles is dynamically incorporated into an MPC framework for mobile robotics, improving efficiency and safety metrics. Other studies encode context implicitly through learned cost maps. For instance, \cite{goel2023semantically} learns dense navigation cost maps in which contextual relationships between objects and places are captured via spatially varying costs. In \cite{10920056}, nonlinear MPC is applied to car-following in autonomous driving, where the speed and position of the preceding vehicle are used as contextual inputs, enabling more precise tracking of the desired inter-vehicle distance.


A popular framework in this direction is Perception-Aware MPC \cite{8593739,bonzanini2024perception,9867264}, where estimates of environmental states are incorporated into the MPC formulation, effectively parametrizing the optimization problem according to the current environment.

Although these approaches demonstrate strong empirical performance, providing closed-loop guarantees in the presence of time-varying contextual signals remains challenging. Naïvely incorporating such signals into costs and constraints may invalidate standard recursive feasibility and stability proofs. While certain guarantees can be established under specific assumptions in perception-aware MPC, deriving constructive procedures that ensure these properties remains difficult \cite{bonzanini2024perception}. The design of appropriate terminal costs and constraints becomes especially nontrivial beyond the linear setting.


Learning-based approaches like neural network control have also emerged as powerful tools for context-aware control because of their inherent ability to digest complex and large amounts of information. Such controllers are frequently designed to depend explicitly on exogenous signals such as environment measurements \cite{zimmerman2014neural,bahabri2025artificial} or fault indicators \cite{yen2000intelligent}. In \cite{huang2020learning} for example, a learning-based switching controller is proposed to ensure reliability of telecommunication services during faults. It is done by feeding a fault detection signal to the controller, facilitating its ability to address such events. Other works use neural networks to map contextual information to control parameters in parameter-varying systems. In \cite{kon2023directlearningparametervaryingfeedforward}, a time-varying scheduling signal is learned from input-output data, and controller gains are mapped through this signal via a neural network. While often effective in practice, the learning-based approaches presented here typically provide limited closed-loop guarantees \cite{zhang2022neural}.

The Performance Boosting (PB) framework introduced in \cite{10633771} aims to bridge this gap by combining the flexibility of neural network controllers with provable closed-loop guarantees in a state-feedback setting. PB control design amounts to unconstrained optimization over state-feedback policies characterized by specific classes of neural networks, that inherently preserve the $\mathcal{L}_p$ stability property of a nonlinear system. In subsequent work \cite{11312988}, this framework was extended to reference tracking, allowing the neural network to incorporate the reference signal as an additional exogenous input while retaining steady-state guarantees. Promising results for fault compensation in drive systems were demonstrated in \cite{kirsch2025resilient}.

This paper extends the PB framework to general contextual signals via an operator decomposition approach, enabling context-enriched control policies that preserve the underlying stability guarantees.


\section{Contributions}
The main contribution of this paper is a structured decomposition for context-enriched PB operators that map process noise and contextual signals to a corrective multi-input control action while preserving $\mathcal L_p$-stability. The proposed factorization isolates a stable disturbance-processing dynamical module from a uniformly bounded context-dependent mixing mechanism, providing an interpretable architecture that can simplify learning-based synthesis and dramatically improve closed-loop performance. Under a tail regularity condition aligned with typical PB disturbance profiles (e.g., transient disturbances), we further establish that this factorization is necessary and sufficient for the associated template class of causal $\mathcal L_p$-stable operators. We also present a numerical experiment illustrating the resulting gains in a representative PB setting.






\section{Notation and Basic Definitions}

Vectors are denoted by lowercase letters, matrices by uppercase letters, and sequences by bold letters.
For a sequence $\mathbf{A}=(A_t)_{t\ge 0}$, we denote by $A_t^{i,j}$ the $(i,j)$-th entry of $A_t$.

\begin{definition}[$\ell^{n\times m}$ and $\ell_p^{n\times m}$ spaces]
Let $\ell^{n\times m}$ be the vector space of all matrix-valued sequences
$\mathbf{A}=(A_t)_{t\ge 0}$ with $A_t\in\mathbb R^{n\times m}$.
For $p\in[1,\infty)$, define
\[
\|\mathbf{A}\|_{p} := \left(\sum_{t=0}^{\infty} \|A_t\|^{p}\right)^{\frac{1}{p}},
\]
and let $\ell_p^{n\times m}:=\{\mathbf{A}\in \ell^{n\times m}:\|\mathbf{A}\|_p<\infty\}$.
Moreover, $\ell_\infty^{n\times m}:=\{\mathbf{A}\in \ell^{n\times m}:\|\mathbf{A}\|_\infty<\infty\}$ where
$\|\mathbf{A}\|_\infty := \sup_{t\ge 0}\|A_t\|$.
\end{definition}


\begin{definition}[Causal operator]
An operator $T:\ell^n\to \ell^m$ is causal if
$T(\mathbf{x}) = (T_0(x_0),\,T_1(x_{0:1}),\,\dots,\,T_t(x_{0:t}),\,\dots)$.
\end{definition}

\begin{definition}[$\mathcal L_p$-stable operator]
An operator $T:\ell^n\to \ell^m$ is said to be $\mathcal L_p$-stable if it is causal and
$T(\mathbf{x})\in \ell_p^m$ for all $\mathbf{x}\in \ell_p^n$. We write $T\in\mathcal L_p$.
\end{definition}
We now introduce the timewise matrix--vector product that will be used throughout the section to
define the proposed factorization.
\begin{definition}[Timewise matrix--vector product]\label{def:timewise-product}
Let $\mathbf{A}=(A_t)_{t\ge 0}\in \ell^{r\times s}$ and $\mathbf{v}=(v_t)_{t\ge 0}\in \ell^{s}$.
Define $\mathbf{A}\boxtimes \mathbf{v}\in \ell^{r}$ by
\[
(\mathbf{A}\boxtimes \mathbf{v})_t := A_t v_t,\qquad t\ge 0.
\]
\end{definition}



\section{Performance Boosting framework}
\label{sec:pb}

We briefly recall the Performance Boosting (PB) framework and motivate the multi-input operator
studied in this paper.

Consider a pre-stabilized discrete-time system, with dynamics: 
\begin{equation}\label{eq:system_state}
    x_t = f_t(x_{0:t-1},u_{0:t-1})+ w_t, ~~~t= 1,2,\ldots\,,
\end{equation}

where $x_t \in \mathbb{R}^n$ is the state vector, $u_t \in \mathbb{R}^m$ is the control input, and $w_t \in \mathbb{R}^n$ stands for unknown process noise. In operator form, system \ref{eq:system_state} is: 
\begin{equation}
\label{eq:operator_form_state}
    \boldsymbol{x} = \mathbf{F}(\boldsymbol{x},\mathbf{u}) + \mathbf{w}\,,
\end{equation}

where $\mathbf{F}:\ell^n\times \ell^m \rightarrow \ell^n$ is the strictly causal operator such that $\mathbf{F}(\mathbf{x},\mathbf{u}) = (0,f_1(x_0,u_0),\ldots,f_t(x_{0:t-1},u_{0:t-1}),\ldots)$.

The main objective of PB is to design the control signal $\mathbf{u}$ such that it enhances the performance of the pre-stabilized system \eqref{eq:operator_form_state} while preserving its steady-state guarantees. To do so, PB relies on an internal-model-control (IMC) principle: a copy of the plant dynamics is used to reconstruct the disturbance signal
\(\mathbf{\widehat w=x -F(x,u)}\) affecting the evolution of the plant. The boosting control action is then parametrized through an operator acting on a this reconstructed signal:
\[
\mathbf{u}= \operator(\mathbf{\widehat w}),
\]
where \(\operator\) is a causal \(\mathcal L_p\)-stable operator to be designed.


A central result of the PB framework is that, under the nominal assumptions of the framework, the search for performance-boosting, stability-preserving controllers can be
reduced to the search for suitable \(\mathcal L_p\)-stable operators \(\operator\). Indeed, the PB architecture provides a parametrization of all and only stability-preserving performance-boosting
controllers in terms of the single operator \(\operator\). This is especially attractive from a learning
and optimization viewpoint: one can optimize over expressive classes of operators while retaining
closed-loop stability by construction. The operator \(\operator\) can be parametrized to minimize any arbitrary loss, while still preserving closed loop guarantees by design.


\begin{comment}
    Consider a pre-stabilized discrete-time system, where a baseline control architecture guarantees the
desired closed-loop stability property, and an additional control input is introduced to improve
transient performance. The key idea behind PB is to parametrize this extra control action through an
operator acting on a disturbance signal reconstructed from the closed loop.

More precisely, PB relies on an internal-model-control (IMC) principle: a copy of the plant dynamics is used to reconstruct the disturbance signal
\(\disturbance\) affecting the evolution of the plant. The boosting control input is then generated as
\[
\mathbf{u}= \operator(\disturbance),
\]
where \(\operator\) is a causal \(\mathcal L_p\)-stable operator to be designed.

A central result of the PB framework is that the search for performance-boosting controllers can be
reduced to the search for suitable \(\mathcal L_p\)-stable operators \(\operator\). In particular,
under the nominal assumptions of the framework, all controllers generated through the PB architecture
with an \(\mathcal L_p\)-stable operator \(\operator\) preserve the \(\mathcal L_p\)-stability of the
closed loop; conversely, every causal controller preserving the same closed-loop \(\mathcal L_p\)
stability can be represented through the PB construction for a suitable operator \(\operator \in \mathcal L_p\).

Therefore, PB provides a parametrization of all and only the stability-preserving performance-boosting
controllers in terms of a single operator \(\operator\). This is especially attractive from a learning
and optimization viewpoint: one can optimize over expressive classes of operators while retaining
closed-loop stability by construction.

\end{comment}

The standard PB formulation uses an \(\mathcal{L}_p\)-stable operator depending only on the reconstructed
disturbance \(\disturbance\), and is in principle able to recover every stability-preserving controller
in this class. However, in many applications it is greatly beneficial to enrich the controller with additional
contextual information. Such context may include state-dependent features, task parameters, operating
conditions, or environment-dependent signals that do not naturally belong to \(\disturbance\) itself and, crucially, do not need to be $ \ell_p$ sequences,
but can considerably simplify the synthesis of high-performance controllers. This motivates the study
of an extended operator of the form
\[
\operator:\ell_p^n\times \ell^m \to \ell_p^r,
\qquad
(\disturbance,\mathbf z)\mapsto \mathbf{u},
\]
where 
\(\mathbf z\in\ell^m\) is a contextual sequence and \(\disturbance\in\ell_p^n\) the disturbance reconstruction.

\begin{comment}
The goal of this paper is to show that such extended operators can be realized through a structured
factorization in which a disturbance-driven \(\mathcal L_p\)-stable dynamical block is combined with a
bounded context-dependent mixer. Moreover, on a broad class of disturbance regimes relevant in control,
this factorization is not conservative.

    
\end{comment}

\section{Main Results}
The main contribution of this paper is to show that extended operators
We consider an \emph{extended} causal operator
\[
\operator:\ell_p^n\times \ell_\infty^m\to \ell_p^r
\]
mapping a disturbance signal \(\disturbance\in\ell_p^n\) and a bounded contextual sequence
\(\mathbf z\in\ell_\infty^m\) to an output \(\operator(\disturbance,\mathbf z)\in\ell_p^r\).
We study a structured factorization that separates
(i) a bounded (``\(\ell_\infty\)'') matrix-valued component depending on \((\disturbance,\mathbf z)\),
and
(ii) a disturbance-driven \(\ell_p\) component depending on \(\disturbance\) only.
The idea is to keep disturbance processing in a causal
\(\mathcal L_p\)-stable block \(\operator_p:\ell_p^n\to\ell_p^s\), and inject context through a bounded
mixer
\[
\operator_\infty:\ell_p^n\times \ell_\infty^m\to \ell_\infty^{r\times s}.
\]
The overall map is obtained via the timewise product.

Let \(\operator_\infty:\ell_p^n\times \ell_\infty^m\to \ell_\infty^{r\times s}\)
and \(\operator_p:\ell_p^n\to \ell_p^{s}\) be causal. Define
\begin{equation}\label{eq:op-decomp}
\operator(\disturbance,\mathbf z)
:= \operator_\infty(\disturbance,\mathbf z)\boxtimes \operator_p(\disturbance).
\end{equation}

This decomposition is naturally motivated by PB architectures: it provides a modular way to inject
context through a bounded mixing mechanism, while keeping the stability-critical disturbance
processing in a separate \(\mathcal L_p\)-stable block. Beyond being sufficient for preserving
\(\ell_p\)-stability, we will show that this factorization is also \emph{non-conservative} on a
PB-relevant disturbance regime: under a mild regularity assumption, typical of transient perturbations
in control, every admissible extended operator admits a representation of the form
\eqref{eq:op-decomp}.

Notice, moreover, that the factorization \ref{eq:op-decomp} generalizes the MAD policies introduced in \cite{furieri2025mad} as well as the rPB policies introduced in \cite{11312988}. Indeed, MAD policies
decompose the control input into an \(\ell_p\)-stable magnitude term and a bounded direction term
depending on the state trajectory. In our framework, this is recovered as the special case \(s=1\),
where \(\operator_p(w)\) generates the scalar magnitude and \(\operator_\infty(w,\mathbf z)\) acts as
a bounded \(r\times 1\) mixer, with \(\mathbf z\) chosen as the state sequence (or a bounded function
thereof). The rPB framework corresponds to the special case where the mixer \(\operator_\infty(w,\mathbf z)\) is diagonal. Our formulation is more general because it allows arbitrary bounded contextual signals,
multi-dimensional disturbance features, and matrix-valued bounded mixing across multiple control
channels.


\subsection{Operator decomposition and \texorpdfstring{$\mathcal L_p$}{Lp}-stability}
We first show that the proposed factorization is always sufficient to preserve
\(\mathcal L_p\)-stability: the timewise product of a bounded mixer and an \(\ell_p\) disturbance
feature sequence remains in \(\ell_p\).
Throughout, fix a vector norm \(\|\cdot\|\) on \(\mathbb R^s\) and the associated induced matrix norm,
so that \(\|A v\|\le \|A\|\,\|v\|\) for all compatible \(A,v\).

\begin{theorem}[Sufficient condition]\label{thm:suff}
Assume \(\operator_\infty(\disturbance,\mathbf z)\in \ell_\infty^{r\times s}\) for all
\(\disturbance\in\ell_p^n\), \(\mathbf z\in\ell_\infty^m\), and
\(\operator_p(\disturbance)\in \ell_p^{s}\) for all \(\disturbance\in\ell_p^n\).
Then \(\operator:\ell_p^n\times \ell_\infty^m\to \ell_p^r\) is well-defined and causal, and
\begin{equation}\label{eq:suff-pointwise}
\|\operator(\disturbance,\mathbf z)\|_p
\le
\|\operator_\infty(\disturbance,\mathbf z)\|_\infty\,
\|\operator_p(\disturbance)\|_p,
\quad \forall(\disturbance,\mathbf z).
\end{equation}
\end{theorem}

\smallskip

\begin{proof}
For every \(t\ge 0\),
\begin{align*}
\|\operator(\disturbance,\mathbf z)_t\|
& =
\|\operator_\infty(\disturbance,\mathbf z)_t\,
\operator_p(\disturbance)_t\| \\
& \le
\|\operator_\infty(\disturbance,\mathbf z)_t\|\,
\|\operator_p(\disturbance)_t\|.
\end{align*}
Therefore,
\[
\sum_{t=0}^\infty \|\operator(\disturbance,\mathbf z)_t\|^p
\le
\Big(\sup_{t\ge 0}\|\operator_\infty(\disturbance,\mathbf z)_t\|\Big)^p
\sum_{t=0}^\infty \|\operator_p(\disturbance)_t\|^p.
\]
Since \(\operator_\infty(\disturbance,\mathbf z)\in\ell_\infty^{r\times s}\) and
\(\operator_p(\disturbance)\in\ell_p^s\), it follows that
\(\operator(\disturbance,\mathbf z)\in\ell_p^r\), and taking the \(p\)-th root yields
\eqref{eq:suff-pointwise}.
\end{proof}

\subsection{A non-conservative factorization on the typical PB regime}

Theorem~\ref{thm:suff} shows that the proposed factorization is always sufficient to preserve
\(\mathcal L_p\)-stability. We now ask when the converse is true, namely when an extended operator
\(\operator\) can be represented in the form \eqref{eq:op-decomp}. This is not automatic on the full
domain \(\ell_p^n\times \ell_\infty^m\). We therefore restrict attention to what we call the
\emph{typical PB regime}, in which disturbances are transient in a uniform weighted sense. On this
regime, we will show that the proposed factorization is non-conservative. The typical PB regime is
characterized by two ingredients: (i) an admissible class of disturbance signals \(\disturbance\),
and (ii) a corresponding regularity property of the operator \(\operator\). We start by introducing
the disturbance class.

\begin{definition}[Weighted \texorpdfstring{$\ell_\infty$}{linfty} envelope space]
\label{def:ellinfty-kappa}
Fix \(p\in[1,\infty)\) and let \(\kappa=(\kappa_t)_{t\ge 0}\) be a positive sequence such that
\begin{equation}\label{eq:kappa-summable}
\sum_{t=0}^{\infty}\kappa_t^{-p}<\infty.
\end{equation}
For \(d\in\mathbb N\), define
\[
\ell_{\infty,\kappa}^d
:=
\left\{
\mathbf v\in \ell^d:\ \|\mathbf v\|_{\kappa,\infty}<\infty
\right\},
\]
with
\[
\|\mathbf v\|_{\kappa,\infty}
:=
\sup_{t\ge 0}\kappa_t\|v_t\|.
\]
\end{definition}

The space \(\ell_{\infty,\kappa}^d\) collects sequences whose magnitude is uniformly dominated by the
time-varying template \(\kappa_t^{-1}\). Indeed, \(\|\mathbf v\|_{\kappa,\infty}<\infty\) is equivalent to
the existence of a constant \(C<\infty\) such that
\[
\|v_t\|\le C\,\kappa_t^{-1}\qquad \forall t\ge 0,
\]
so \(\kappa_t^{-1}\) plays the role of a prescribed decay (or envelope) profile.
The summability condition \eqref{eq:kappa-summable} ensures that this template is \(\ell_p\)-summable,
and therefore any \(\mathbf v\in\ell_{\infty,\kappa}^d\) automatically belongs to \(\ell_p^d\), with the
explicit bound
\begin{equation}\label{eq:kappa-embed-inline}
\|\mathbf v\|_p
\le
\left(\sum_{t=0}^\infty \kappa_t^{-p}\right)^{\!\frac1p}
\|\mathbf v\|_{\kappa,\infty}.
\end{equation}

The admissible class \(\ell_{\infty,\kappa}^n\) is also broad: it contains all finitely supported
sequences (since \(\kappa_t>0\) for all \(t\)) and is therefore dense in \(\ell_p^n\) for every
\(p\in[1,\infty)\). Moreover, \(\kappa\) should be interpreted as a \emph{worst-case} template:
if a disturbance decays faster than \(\kappa_t^{-1}\), then it still belongs to
\(\ell_{\infty,\kappa}^n\). In this sense, faster transients make the assumption easier to satisfy,
and one may deliberately choose a \emph{slower} (hence less restrictive) envelope, for instance a
polynomial one
$\kappa_t=(t+1)^\alpha,\: \alpha>\frac{1}{p},$
to capture a broad admissible disturbance regime.

\smallskip

We now introduce the second ingredient of the typical PB regime, namely a mild regularity condition on
the extended operator \(\operator\) over the admissible disturbance class.

\begin{assumption}[Envelope-preserving operator]\label{ass:env}
Let \(\operator\) be causal on \(\ell_{\infty,\kappa}^n\times \ell_\infty^m\) with values in \(\ell_p^r\).
Assume that for every \(\disturbance\in\ell_{\infty,\kappa}^n\) and
\(\mathbf z\in\ell_\infty^m\),
\begin{equation}\label{eq:env}
\sup_{t\ge 0}\kappa_t\,\|\operator(\disturbance,\mathbf z)_t\| < \infty.
\end{equation}
\end{assumption}

Assumption~\ref{ass:env} requires the operator \(\operator\) to preserve the admissible envelope class:
whenever the disturbance \(\disturbance\) is dominated by the template \(\kappa_t^{-1}\), the output
\(\operator(\disturbance,\mathbf z)\) must remain dominated by the same template, up to a multiplicative
constant that may depend on \((\disturbance,\mathbf z)\) but is uniform in time. This is exactly the
additional regularity that makes a bounded-mixer factorization possible on the typical PB regime, as
shown by the next result. Moreover, this assumption is not merely abstract: it is satisfied by a broad
family of dynamical operators commonly used in modern control applications, as we shall see later on.

Under Assumption~\ref{ass:env}, we can establish a partial converse to
Theorem~\ref{thm:suff}: within the typical PB regime, the bounded-mixer factorization is not only
sufficient for \(\mathcal L_p\)-stability, but it is also \emph{necessary} for operators satisfying the
same uniform envelope regularity. Overall, this shows that on the typical PB regime the proposed
factorization is non-conservative.

\begin{theorem}[Necessary and sufficient factorization]
\label{thm:equiv}
Fix \(p\in[1,\infty)\), a weight sequence \(\kappa\) satisfying \eqref{eq:kappa-summable}, and \(s\ge 1\).
Let \(\operator\) be causal on \(\ell_{\infty,\kappa}^n\times \ell_\infty^m\). Then the following are equivalent:
\begin{enumerate}
\item[\textup{(i)}] \textbf{Uniform envelope.}
For every \(\disturbance\in\ell_{\infty,\kappa}^n\) and \(\mathbf z\in\ell_\infty^m\),
\begin{equation}\label{eq:env-iff}
\sup_{t\ge 0}\kappa_t\,\|\operator(\disturbance,\mathbf z)_t\|<\infty.
\end{equation}

\item[\textup{(ii)}] \textbf{Envelope-compatible factorization.}
There exist causal operators
\[
\operator_p:\ell_{\infty,\kappa}^n\to \ell_{\infty,\kappa}^s,
\qquad
\operator_\infty:\ell_{\infty,\kappa}^n\times \ell_\infty^m
\to \ell_\infty^{r\times s},
\]
such that
\begin{equation}\label{eq:env-factor}
\operator(\disturbance,\mathbf z)
=
\operator_\infty(\disturbance,\mathbf z)\boxtimes \operator_p(\disturbance),
\qquad
\forall(\disturbance,\mathbf z).
\end{equation}
\end{enumerate}
\end{theorem}

\begin{proof}
\emph{(i)$\Rightarrow$(ii).}
Assume \eqref{eq:env-iff}. It is enough to exhibit one admissible factorization.
Define the causal map
\begin{equation}\label{eq:vp-canonical}
\operator_p(\disturbance)_t := \kappa_t^{-1}\,\mathbf 1_s,\qquad t\ge 0,
\end{equation}
and define \(\operator_\infty\) by
\begin{equation}\label{eq:Minfty-canonical}
\big(\operator_\infty(\disturbance,\mathbf z)_t\big)^{i,j}
:=
\begin{cases}
\kappa_t\,\operator(\disturbance,\mathbf z)_t^{i}, & j=1,\\
0, & j\neq 1.
\end{cases}
\end{equation}
Then \eqref{eq:Minfty-canonical}--\eqref{eq:vp-canonical} give
\[
\operator_\infty(\disturbance,\mathbf z)_t\,
\operator_p(\disturbance)_t
=
\operator(\disturbance,\mathbf z)_t
\]
for all \(t\), hence \eqref{eq:env-factor} holds.

Moreover, \(\operator_p(\disturbance)\in \ell_{\infty,\kappa}^s\) since
\[
\|\operator_p(\disturbance)\|_{\kappa,\infty}
=
\sup_{t\ge 0}\kappa_t\|\kappa_t^{-1}\mathbf 1_s\|
=
\|\mathbf 1_s\|<\infty.
\]
Finally, \eqref{eq:env-iff} implies
\(\operator_\infty(\disturbance,\mathbf z)\in \ell_\infty^{r\times s}\).
Indeed, \(\operator_\infty(\disturbance,\mathbf z)_t\) has only one nonzero column equal to
\(\kappa_t\,\operator(\disturbance,\mathbf z)_t\in\mathbb R^r\), hence there exists \(c>0\)
(depending only on the chosen norm on \(\mathbb R^s\)) such that
\[
\|\operator_\infty(\disturbance,\mathbf z)_t\|
\le
c\,\kappa_t\|\operator(\disturbance,\mathbf z)_t\|,
\qquad \forall t\ge 0.
\]
Taking the supremum over \(t\) yields
\(\|\operator_\infty(\disturbance,\mathbf z)\|_\infty<\infty\).

\smallskip
\emph{(ii)$\Rightarrow$(i).}
Assume \eqref{eq:env-factor}. For each \(t\ge 0\),
\begin{align*}
\kappa_t\|\operator(\disturbance,\mathbf z)_t\|
 & =
\kappa_t\|\operator_\infty(\disturbance,\mathbf z)_t
\,\operator_p(\disturbance)_t\| \\
& \le
\|\operator_\infty(\disturbance,\mathbf z)\|_\infty\,
\kappa_t\|\operator_p(\disturbance)_t\|.
\end{align*}
Taking the supremum over \(t\) gives
\[
\sup_{t\ge 0}\kappa_t\|\operator(\disturbance,\mathbf z)_t\|
\le
\|\operator_\infty(\disturbance,\mathbf z)\|_\infty\,
\|\operator_p(\disturbance)\|_{\kappa,\infty}
<\infty,
\]
since \(\operator_p(\disturbance)\in \ell_{\infty,\kappa}^s\). This proves \eqref{eq:env-iff}.
\end{proof}

\begin{remark}[Design viewpoint and ``typical PB regime'']\label{rem:typical-pb}
Theorem~\ref{thm:equiv} shows that, on the admissible regime, the proposed factorization is
non-conservative: every envelope-regular extended operator admits a representation of the form
\(\operator=\operator_\infty\boxtimes \operator_p\).
From a synthesis viewpoint, however, one typically proceeds in the opposite direction and starts from
item~\textup{(ii)} of Theorem~\ref{thm:equiv}, i.e., one designs the two factors
\(\operator_p\) and \(\operator_\infty\) directly.
Theorem~\ref{thm:suff} then guarantees that the resulting operator is \(\mathcal L_p\)-stable by
construction, which is precisely the property required in PB to ensure overall closed-loop stability.
Theorem~\ref{thm:equiv} further clarifies that, for the disturbance regimes typically encountered in
control applications and for the broad class of operators satisfying Assumption~\ref{ass:env}, this
architectural choice is not conservative.

In particular, operators realized by asymptotically stable (fading-memory) systems with exponential
forgetting preserve rich classes of admissible envelopes, including all polynomial choices
\(\kappa_t=(t+1)^\alpha\).
As a consequence, the admissible class \(\ell_{\infty,\kappa}^n\) can be chosen to accommodate a large
subset of \(\ell_p\) disturbances encountered in practice, including finite-support transients and
disturbances with polynomially decaying tails. We formalize this point in the next result.
\end{remark}

\begin{lemma}[Fading memory preserves polynomial envelopes]
\label{lem:fading-poly}
Fix \(\alpha>0\) and let \(\kappa_t=(t+1)^\alpha\).
Assume that a causal operator \(M:\ell^n\to\ell^r\) satisfies the fading-memory estimate:
there exist constants \(a,b\ge 0\) and \(\rho\in(0,1)\) such that for every input \(u\) and every \(t\ge 0\),
\begin{equation}\label{eq:fading}
\|M(u)_t\|
\le
a\rho^t + b\sum_{k=0}^{t}\rho^{t-k}\|u_k\|.
\end{equation}
Then \(M\) maps \(\ell_{\infty,\kappa}^n\) into \(\ell_{\infty,\kappa}^r\) and
\begin{equation}\label{eq:fading-poly-bound}
\|M(u)\|_{\kappa,\infty}
\le
a\,\sup_{t\ge 0}(t+1)^\alpha\rho^t
+
b\,C_{\alpha,\rho}\,\|u\|_{\kappa,\infty},
\end{equation}
where
\[
C_{\alpha,\rho}
:=
\sum_{j=0}^\infty \rho^{j}(j+1)^\alpha
<\infty.
\]
In particular, polynomial envelopes of order \(\alpha\) are preserved.
\end{lemma}

\begin{proof}
Let \(u\in\ell_{\infty,\kappa}^n\), so
\(\|u_k\|\le \|u\|_{\kappa,\infty}(k+1)^{-\alpha}\).
Using \eqref{eq:fading} and the change of variables \(j=t-k\),
\begin{align*}
(t+1)^\alpha\|M(u)_t\|
&\le
a(t+1)^\alpha\rho^t \\
&\quad
+b\|u\|_{\kappa,\infty}
\sum_{k=0}^{t}\rho^{t-k}\frac{(t+1)^\alpha}{(k+1)^\alpha} \\
&=
a(t+1)^\alpha\rho^t \\
&\quad
+b\|u\|_{\kappa,\infty}
\sum_{j=0}^{t}\rho^{j}
\frac{(t+1)^\alpha}{(t-j+1)^\alpha}.
\end{align*}
Since \(t-j+1\ge 1\), we have
\[
\frac{(t+1)^\alpha}{(t-j+1)^\alpha}\le (j+1)^\alpha,
\]
hence
\[
(t+1)^\alpha\|M(u)_t\|
\le
a(t+1)^\alpha\rho^t
+
b\|u\|_{\kappa,\infty}
\sum_{j=0}^{\infty}\rho^{j}(j+1)^\alpha.
\]
Taking the supremum over \(t\) yields \eqref{eq:fading-poly-bound}. The series defining
\(C_{\alpha,\rho}\) converges because exponential decay dominates polynomial growth.
\end{proof}

The previous lemma is particularly relevant in practice because the envelope-preserving property is
satisfied by a broad class of state-of-the-art dynamical operators used in modern control-oriented learning architectures.
In particular, this is the case for stable structured state-space models such as the Recurrent Equilibrium Unit \cite{orvieto2023resurrecting} and
the L2RU architecture \cite{11312999}, whose internal linear dynamics are asymptotically stable. As a result, the
induced input-output map exhibits fading-memory behavior, i.e., exponential forgetting of past inputs,
and the model preserves all polynomial envelope classes, and more generally, a large family of
moderate-growth subexponential envelopes. Similarly, recurrent equilibrium networks (RENs) \cite{revay2021recurrent} are designed
to be contracting and robust, which places them within the same fading-memory regime. Hence, under their
standard stability assumptions, both stable SSMs and RENs satisfy the type of envelope regularity
required here.
\section{Numerical Experiment: Contextual Gate Navigation}
\label{sec:exp-gate}

To illustrate the practical role of contextual multi-input PB, we consider a planar robot that must
navigate through a thin wall at \(x_{\mathrm{wall}}\) by passing through a gate whose centre
\(g_t\) switches randomly over time and freezes before the robot reaches the wall.
The key challenge is that the gate position is unknown at the start of the episode and can change
during the approach, so the controller must read the context to know where to aim and when it
is safe to commit.

\subsection{Pre-stabilized Dynamics and PB Disturbance}
\label{subsec:gate-dynamics}

Let
\begin{equation}
x_t=
\begin{bmatrix}
p_t^\top & v_t^\top
\end{bmatrix}^\top
\in\mathbb{R}^4,
\qquad
u_t\in\mathbb{R}^2,
\end{equation}
where \(p_t=(p_{x,t},p_{y,t})^\top\) and \(v_t=(v_{x,t},v_{y,t})^\top\) denote position and
velocity, and \(u_t\) is the PB corrective input. The robot follows a pre-stabilized double integrator:
\begin{align}
p_{t+1} &= p_t + \Delta t\, v_t, \label{eq:gate-nom-p}\\
v_{t+1} &= v_t + \Delta t\big(-K_p p_t - K_d v_t + u_t\big),
\label{eq:gate-nom-v}
\end{align}
with fixed PD gains \(K_p,K_d>0\) and timestep \(\Delta t=0.05\,\mathrm{s}\). We denote by
\(f_{\mathrm{nom}}\) the nominal closed-loop map, so that the true rollout satisfies
\begin{equation}
x_0 = w_0,
\qquad
x_{t+1}=f_{\mathrm{nom}}(x_t,u_t)+w_{t+1},
\quad t\ge 0.
\label{eq:gate-rollout}
\end{equation}

The additive disturbance \(w_t\in\mathbb{R}^4\) combines small persistent state noise with a
small number of intermittent wind gusts. Specifically, at each step position and velocity
components receive zero-mean Gaussian perturbations with standard deviations
\(\sigma_{\mathrm{pos}}\) and \(\sigma_{\mathrm{vel}}\), respectively. Additionally,
\(n_g\sim\mathrm{Uniform}\{n_{\min},n_{\max}\}\) gusts occur per episode, each lasting
\(d_g\) steps; during a gust, a lateral velocity impulse in the range
\([\delta_{\min},\delta_{\max}]\) is added to \(v_{y,t}\), clipped to prevent excessive
deviation.

Following the PB formulation, the disturbance is reconstructed from one-step nominal prediction errors:
\begin{align}
w_0 &= x_0, \label{eq:gate-w0}\\
w_t &= x_t-f_{\mathrm{nom}}(x_{t-1},u_{t-1}),
\qquad t\ge 1.
\label{eq:gate-wt}
\end{align}
Under \eqref{eq:gate-rollout} this reconstruction is exact, so the same signal \(\mathbf{w}\) both
perturbs the plant and is fed to the PB operator.

\subsection{Moving Gate, Context Signal, and Lifting}
\label{subsec:gate-context}

\paragraph{Gate schedule.}
The gate centre \(g_t\in[-g_{\max},g_{\max}]\) follows a piecewise-constant random
schedule. Starting from \(g_0=0\), the gate holds each level for a dwell time drawn uniformly
from \(\{d_{\min},\dots,d_{\max}\}\) steps, then jumps to a new level. The schedule is
frozen at a fixed step
\begin{equation}
t_{\mathrm{freeze}}
=
t_{\mathrm{wall}}-n_{\mathrm{settle}},
\label{eq:freeze-step}
\end{equation}
where \(t_{\mathrm{wall}}\) is the expected wall-crossing step under nominal zero-input dynamics
from the midpoint initial condition, and \(n_{\mathrm{settle}}\ge 1\) is a small settling margin.
For \(t\ge t_{\mathrm{freeze}}\) the gate is constant at \(g_{t_{\mathrm{freeze}}-1}\).
The value of \(t_{\mathrm{wall}}\) is precomputed once by simulating the nominal plant forward
until \(p_{x,t}\le x_{\mathrm{wall}}\), and is the same for all episodes.

To stress the controller, \(30\%\) of training and test episodes include an \emph{adversarial
late switch}: the gate jumps to a new level at a step drawn uniformly in
\([t_{\mathrm{freeze}}-d_{\min},\, t_{\mathrm{freeze}})\), giving the robot very little time to
adapt before the gate freezes.

\paragraph{Context signal.}
At each step \(t\) the controller observes a time-varying context vector \(z_t\in\mathbb{R}^{n_z}\).
We study two variants:
\begin{itemize}
  \item \textbf{Full context} (\(n_z=11\)): comprises the current gate centre \(g_t\), its finite
    difference \(\dot g_t\), a fast and a slow exponential moving average of \(g_t\), the
    signed gate error \(e_t:=p_{y,t}-g_t\), the relative wall distance
    \(\delta_t:=p_{x,t}-x_{\mathrm{wall}}\), the goal offsets \((d_{x,t},d_{y,t})\), an
    approach weight \(\alpha_t\), the normalised switch age \(\sigma_t\), and the
    normalised time-to-wall \(\tau_t\).

  \item \textbf{Minimal context} (\(n_z=3\)): contains only the three features that are
    directly observable without privileged knowledge of the freeze schedule:
    \begin{equation}
      z_t^{\min}
      =
      \begin{bmatrix}
        e_t/y_{\mathrm{sc}} & \alpha_t & \sigma_t/T
      \end{bmatrix}^\top,
      \label{eq:minimal-context}
    \end{equation}
    where \(e_t = p_{y,t}-g_t\) is the gate error (where to aim),
    \(\alpha_t = \exp\!\bigl(-\tfrac{1}{2}(\delta_t/\sigma_{\alpha})^2\bigr)\) is a
    Gaussian proximity weight centred at the wall (spatial commitment cue), and
    \(\sigma_t\) is the number of steps since the last gate switch (stability cue:
    large \(\sigma_t\) means the gate has been settled long enough to commit).
\end{itemize}
The three features are directly observable without privileged knowledge of the freeze
schedule \eqref{eq:freeze-step}: a rising \(\sigma_t\) near the wall
(\(\alpha_t\approx 1\)) indicates that the gate has been stable long enough to aim,
providing the information necessary to navigate the trade-off between early commitment
and waiting for the gate to settle.

\paragraph{Context lifting.}
Rather than feeding \(z_t\) directly into \(\mathcal{M}_p\), we first lift it through a
small causal network \(\psi:\mathbb{R}^{n_z}\to\mathbb{R}^{d_\psi}\) and augment the
disturbance input:
\begin{equation}
\bar{w}_t = \bigl[w_t^\top,\;\psi(z_t)^\top\bigr]^\top \in \mathbb{R}^{n_x+d_\psi}.
\label{eq:lifting}
\end{equation}
The operator \(\mathcal{M}_p\) then maps \(\bar{\mathbf{w}}\) to features, while
\(\mathcal{M}_b\) sees the original \((w,z)\).

\paragraph{Factorized PB operator.}
The full PB correction uses the factorized form
\begin{equation}
u_t = \mathcal{M}_b(w,z)_t\,v_t,
\qquad
v_t = \mathcal{M}_p(\bar{w})_t,
\label{eq:gate-factor}
\end{equation}
where \(\mathcal{M}_p\) is a stable SSM (LRU parametrization, \(L\) layers,
\(d_{\mathrm{model}}\) hidden units) that processes the lifted disturbance sequence, and
\(\mathcal{M}_b\) is a bounded MLP with spectral-norm layers and soft-sign output that produces
a gain matrix \(A_t\in\mathbb{R}^{2\times s}\).

Since this architecture falls within the typical PB regime of Section~\ref{sec:pb},
Theorem~\ref{thm:equiv} guarantees that the factorization is non-conservative.

\paragraph{Ablation: \(\mathcal{M}_p\)-only.}
We also evaluate a baseline in which \(\mathcal{M}_b\) is removed and the output of the SSM
is used directly as the control correction. To ensure a fair comparison, the SSM hidden
dimension is enlarged so that its total parameter count matches the combined parameter budget
of \(\mathcal{M}_p + \mathcal{M}_b\) in the factorized variant.

\paragraph{Dataset.}
The training dataset consists of \(N\) on-policy rollouts generated online. Each rollout is
defined by an initial condition \(x_0^{(k)}\), a disturbance sequence \(\mathbf{w}^{(k)}\),
and a gate schedule \(g_{0:T}^{(k)}\). Training uses paired episodes: each initial
condition is paired with two independent gate schedules, encouraging the controller to adapt
to different gate configurations from the same start state.

\subsection{Training Objective}
\label{subsec:gate-training}

For a rollout of horizon \(T\), the trajectory cost is
\begin{align}
\mathcal{J}^{(k)}
=\;&
\lambda_{\mathrm{goal}}\,\|p_T^{(k)}\|^2
+
\frac{\lambda_{\mathrm{wall}}}{T}\sum_{t=0}^{T-1}
  \alpha_t^{(k)}\,e_t^{(k)\,2}
\notag\\
&+
\frac{\lambda_{\mathrm{gate}}}{T}\sum_{t=0}^{T-1}
  \kappa_t^{(k)}\,
  \mathrm{ReLU}\!\bigl(|e_t^{(k)}|-h\bigr)
+
\frac{\lambda_{\mathrm{cor}}}{T}\sum_{t=0}^{T-1}
  \phi_{\mathrm{cor}}\!\bigl(p_{y,t}^{(k)}\bigr),
\label{eq:gate-loss}
\end{align}
where \(h\) is the gate half-width, \(\kappa_t^{(k)}\) is a sharpness weight concentrated
near the wall, and \(\phi_{\mathrm{cor}}\) is a corridor-limit penalty. The first term
drives the robot to the origin; the second penalises lateral deviation at the wall weighted
by the approach proximity \(\alpha_t\); the third adds a hard margin once the gate error
exceeds the half-width; the fourth enforces the corridor bounds.
The training objective minimises the empirical average \(\mathcal{J}_{\mathrm{train}}=\frac{1}{N}\sum_k\mathcal{J}^{(k)}\) over \(\theta\).

\subsection{Evaluation Metrics}
\label{subsec:gate-metrics}

We report the following metrics on a held-out test set:
\begin{itemize}
\item \textbf{Gate success rate}: fraction of episodes where the robot crosses the wall within
  the gate opening, i.e., \(|e_{t^\star}^{(k)}|\le h\), where \(t^\star\) is the first step
  with \(p_{x,t^\star}\le x_{\mathrm{wall}}\).
\item \textbf{Mean cross error}: average \(|e_{t^\star}^{(k)}|\) over all test episodes.
\item \textbf{Adversarial success rate}: gate success rate restricted to adversarial episodes
  (switch age at crossing \(\le d_{\min}\)), measuring robustness to late gate changes.
\end{itemize}

We compare four variants: (i)~\textbf{Nominal}: PD baseline, no PB correction;
(ii)~\textbf{Disturbance only}: \(\mathcal{M}_p\)-only SSM with no context (\(z_t=0\));
(iii)~\textbf{Factorized + context}: full \(\mathcal{M}_b\boxtimes\mathcal{M}_p\) with
lifting; and (iv)~\textbf{\(\mathcal{M}_p\)-only + context}: parameter-matched SSM baseline.
Results are reported for both the full (\(n_z=11\)) and minimal (\(n_z=3\)) context variants.

\begin{table}[t]
\centering
\caption{Gate navigation results on the held-out test set (mean over 3 seeds).}
\label{tab:gate-results}
\begin{tabular}{lcccc}
\hline
\textbf{Method} & \textbf{Context} & \textbf{Succ.\ (\%)} & \textbf{Cross err.} & \textbf{Adv.\ succ.\ (\%)} \\
\hline
Nominal            & --      & -- & -- & -- \\
Disturbance only   & --      & -- & -- & -- \\
Factorized + lift  & Full    & -- & -- & -- \\
Factorized + lift  & Minimal & -- & -- & -- \\
\(\mathcal{M}_p\)-only + lift & Full    & -- & -- & -- \\
\(\mathcal{M}_p\)-only + lift & Minimal & -- & -- & -- \\
\hline
\end{tabular}
\end{table}

\subsection{Figures}
\label{subsec:gate-figs}

\begin{figure}[t]
\centering
\fbox{\parbox[c][3.6cm][c]{0.95\columnwidth}{\centering
Placeholder: success rate and mean cross error bars\\
for all variants, full vs.\ minimal context.}}
\caption{Gate success rate and mean cross error for all variants. The factorized controller with
context outperforms the disturbance-only and \(\mathcal{M}_p\)-only baselines.}
\label{fig:gate_summary}
\end{figure}

\begin{figure}[t]
\centering
\fbox{\parbox[c][3.6cm][c]{0.95\columnwidth}{\centering
Placeholder: adversarial switching panel\\
(success rate bar, cross-error violin, hard trajectories).}}
\caption{Robustness to adversarial late gate switches. Left: success rates split by adversarial
vs.\ stable episodes. Centre: distribution of cross errors for adversarial episodes. Right:
top-down trajectories for the hardest adversarial cases.}
\label{fig:gate_adversarial}
\end{figure}

\begin{figure}[t]
\centering
\fbox{\parbox[c][3.6cm][c]{0.95\columnwidth}{\centering
Placeholder: waiting-behaviour plot\\
(velocity profile near wall, aligned to freeze step).}}
\caption{Learned waiting behaviour. The contextual PB controller decelerates before
\(t_{\mathrm{freeze}}\) and commits once the gate has settled, exploiting the switch-age cue.}
\label{fig:gate_waiting}
\end{figure}


\section{Conclusions}

We proposed a structured decomposition for context-enriched Performance Boosting operators, in which a
disturbance-driven \(\mathcal L_p\)-stable dynamical module is combined with a bounded
context-dependent mixer. This factorization yields a modular and interpretable architecture for
multi-input boosting control, while preserving the \(\mathcal L_p\)-stability guarantees required by
the PB framework.

Our main theoretical result shows that, on a broad admissible disturbance regime characterized by a
uniform envelope condition, the proposed factorization is not only sufficient but also
non-conservative: every envelope-regular extended operator admits a representation of the proposed
form. We further discussed how this regime naturally covers many transient disturbances encountered in
control practice, and how fading-memory dynamical models, including stable SSM-based architectures,
satisfy the corresponding envelope-preserving property.

The numerical experiment illustrates that enriching the PB operator with contextual information can
improve the allocation of corrective actions and lead to better control performance in representative
multi-input scenarios. Future work will investigate richer contextual architectures, tighter classes of
admissible envelopes, and broader applications of the proposed decomposition in learning-based control.


\bibliographystyle{ieeetr}
\bibliography{bib}

\end{document}